% main.tex -- Preparatory textbook for reading <AUTHORS>,
%             "<TITLE>"
%
% Build:  sh tex/check.sh   (never an ad-hoc latexmk + grep; see CLAUDE.md)
% ASCII only under tex/.
\documentclass[11pt,oneside]{book}

\usepackage{preamble}
\usepackage{notation}

% The title block carries the paper's full reference, so that the book states
% on its face what it is preparatory FOR.  Fill it from spec.md, "Target
% paper" -- the same string, not a paraphrase.
\title{\textbf{<Title of the Paper>}\\[0.4em]
       \large A preparatory textbook\\[1.2em]
       \normalsize
       <Authors>,\\
       \emph{<full title>},\\
       <journal / arXiv identifier>, <pages>}
\author{}
\date{}

\begin{document}
\maketitle

\frontmatter
\tableofcontents

\mainmatter

% Front matter proper (Phase 6): what the book is, how the two Parts relate,
% the reading order and why, and the legend for the verification markers.
% \include{preface}

\part{Preparatory chapters}
% Chapters are added ONE AT A TIME, each after approval (CLAUDE.md, Phase 4).
% Uncomment as they are written; the order is the one fixed in spec.md,
% "Chapter plan".
% \include{chapters/ch01-<slug>}
% \include{chapters/ch02-<slug>}

\part{Reading the paper}
% Phase 5.  The order is NOT necessarily the paper's; see spec.md,
% "Phase 5 order".
% \include{paper/sec01-<slug>}

% \part{Appendices}
\appendix
% Reader-facing tables that lived only in the management documents until
% Phase 6 -- a notation correspondence table across several papers, the list
% of misprints found in the source, and so on.
% \include{appendix-notation}
% \include{appendix-misprints}

\backmatter
% \nocite{*} keeps bibtex from failing on a book with no \cite yet.  DELETE IT
% once the first chapter cites anything: it would otherwise print every entry
% of refs.bib, including ones the book never uses.
\nocite{*}
\bibliographystyle{alpha}
\bibliography{refs}

\end{document}
